% Este fuente se publica en dos dependencias causales mediante los
% selectores definidos por los wrappers de integración. Ambos tramos cubren
% conjuntamente el archivo completo, sin solape ni omisión.
\ifdefined\HMTGMCOperatorialAPP
\section[Rango matricial de APP]{Rango matricial de APP: clasificación exacta
del álgebra suma--producto e intervalo libre}
\label{p12-83-c22-s1:gmc:chap:operatorios}

Los cuadrados qutrit del \cref{p12-83-c22-s3:gmc:chap:qms} proporcionan una entrada
operatorial. Existe una
segunda entrada, más interna a APP: las dos particiones inducidas por suma y
producto generan proyecciones ortogonales no conmutativas. El lenguaje de
sistemas de operadores permite estudiarlas sin forzar una cuadrícula.

\subsection{Par de esperanzas condicionales asociado a las particiones aditiva y multiplicativa de APP}

Sea
\[
 \mathcal H_d=\ell^2(\{1,\ldots,9\}^d),
\]
y definamos
\[
 \Sigma_d(x)=\operatorname{dr}(x_1+\cdots+x_d),
 \qquad
 \Pi_d(x)=\operatorname{dr}(x_1\cdots x_d).
 \tag{GMC-6.1}
\]
Sean $P_{\Sigma,d}$ y $P_{\Pi,d}$ las proyecciones ortogonales sobre las
funciones constantes en las fibras de $\Sigma_d$ y $\Pi_d$. El puente
precedente demuestra
\[
 \|[P_{\Sigma,2},P_{\Pi,2}]\|=\frac{\sqrt2}{3},
 \qquad
 \|[P_{\Sigma,3},P_{\Pi,3}]\|=\frac{\sqrt3}{4},
 \tag{GMC-6.2}
\]
y
\[
 \operatorname{Ran}P_{\Sigma,d}\cap
 \operatorname{Ran}P_{\Pi,d}=\mathbb C\mathbf1
 \qquad(2\le d\le15).
 \tag{GMC-6.3}
\]
APP contiene, por tanto, un paquete operatorial no conmutativo antes de
cualquier denominación QLS o QMS.

Definimos
\[
 \mathcal S_d=\Span\{I,P_{\Sigma,d},P_{\Pi,d}\},
 \qquad
 \mathfrak C_d^{\rm APP}=C^*(P_{\Sigma,d},P_{\Pi,d}).
 \tag{GMC-6.4}
\]
$\mathcal S_d$ es un sistema de operadores y su rango matricial es
\[
 \mathcal W_s(\mathrm{APP}_d)=
 \left\{(\Phi(P_{\Sigma,d}),\Phi(P_{\Pi,d})):
 \Phi:\mathfrak C_d^{\rm APP}\to M_s\ \text{UCP}\right\}.
 \tag{GMC-6.5}
\]
La torre $\bigsqcup_s\mathcal W_s$ es matricialmente convexa. A diferencia de
una etiqueta general de «cuanticidad», \textup{(GMC-6.5)} permite calcular extremos,
compresiones, dilataciones e invariantes.

\subsection{Clasificación del álgebra suma--producto en profundidad 2}

Pongamos $P=P_{\Sigma,2}$ y $Q=P_{\Pi,2}$ sobre $\mathcal H_2$, de dimensión
81. Las nueve fibras de $\Sigma_2$ tienen tamaño nueve. Las fibras de
$\Pi_2$ tienen tamaños
\[
 6,6,12,6,6,12,6,6,21.
 \tag{GMC-6.6}
\]
Sean $u_a$ y $v_b$ los vectores indicadores normalizados de esas fibras, y
$C=(\langle u_a,v_b\rangle)_{a,b=1}^9$ su matriz de Gram cruzada. El recuento
exacto de intersecciones produce
\[
 \chi_{CC^*}(t)
 =t^5(t-1)\left(t-\frac57\right)
       \left(t-\frac13\right)^2.
 \tag{GMC-6.7}
\]
Además,
\[
\begin{array}{c|c}
\text{subespacio común} & \text{dimensión}\\
\hline
\ker P\cap\ker Q&64\\
\operatorname{Ran}P\cap\operatorname{Ran}Q&1\\
\operatorname{Ran}P\cap\ker Q&5\\
\ker P\cap\operatorname{Ran}Q&5
\end{array}
\tag{GMC-6.8}
\]

\begin{theorem}[Álgebra suma--producto de APP en $d=2$]
\label{p12-83-c22-s1:gmc:thm:algebra-app-d2}
Se tiene un isomorfismo de $C^*$-álgebras finito-dimensionales
\[
 \boxed{\ 
 \mathfrak C_2^{\rm APP}
 \cong\mathbb C^4\oplus M_2(\mathbb C)\oplus M_2(\mathbb C).
 \ }
 \tag{GMC-6.9}
\]
En particular, $\dim\mathfrak C_2^{\rm APP}=12$.
\end{theorem}

\begin{proof}
La descomposición canónica de Halmos para dos proyecciones separa los cuatro
subespacios de \textup{(GMC-6.8)} y un sector genérico para cada ángulo principal
$0<\lambda<1$. La ecuación \textup{(GMC-6.7)} muestra que los únicos valores
genéricos distintos son $5/7$ y $1/3$. En cada uno, el par es unitariamente
equivalente a
\[
 P_0=\begin{pmatrix}1&0\\0&0\end{pmatrix},
 \qquad
 Q_\lambda=
 \begin{pmatrix}
  \lambda&\sqrt{\lambda(1-\lambda)}\\
  \sqrt{\lambda(1-\lambda)}&1-\lambda
 \end{pmatrix},
 \tag{GMC-6.10}
\]
que genera $M_2$. La multiplicidad de $\lambda=1/3$ amplifica la misma
representación pero no crea otro sumando simple. Los cuatro sectores
simultáneamente $0/1$ producen $\mathbb C^4$.
\end{proof}

En concreto,
\[
 Q_{5/7}=
 \begin{pmatrix}5/7&\sqrt{10}/7\\\sqrt{10}/7&2/7\end{pmatrix},
 \qquad
 Q_{1/3}=
 \begin{pmatrix}1/3&\sqrt2/3\\\sqrt2/3&2/3\end{pmatrix}.
 \tag{GMC-6.11}
\]
La norma del conmutador es el máximo de
$\sqrt{\lambda(1-\lambda)}$ y recupera $\sqrt2/3$.

\begin{corollary}[Parametrización completa del rango matricial]
\label{p12-83-c22-s1:gmc:cor:rango-app-d2}
Un par $(X,Y)$ pertenece a $\mathcal W_s(\mathrm{APP}_2)$ si y sólo si
existen efectos
$E_{00},E_{11},E_{10},E_{01}\succeq0$ y mapas CP
$\Psi_{5/7},\Psi_{1/3}:M_2\to M_s$ tales que
\[
 E_{00}+E_{11}+E_{10}+E_{01}
 +\Psi_{5/7}(I)+\Psi_{1/3}(I)=I_s,
 \tag{GMC-6.12}
\]
\[
 X=E_{11}+E_{10}
 +\Psi_{5/7}(P_0)+\Psi_{1/3}(P_0),
 \tag{GMC-6.13}
\]
\[
 Y=E_{11}+E_{01}
 +\Psi_{5/7}(Q_{5/7})+\Psi_{1/3}(Q_{1/3}).
 \tag{GMC-6.14}
\]
\end{corollary}

\begin{proof}
Todo mapa CP desde una suma directa se descompone en mapas CP sobre sus
sumandos. Los cuatro sumandos escalares producen los efectos $E_{ab}$; los
dos sumandos matriciales producen $\Psi_\lambda$. La unitalidad es
\textup{(GMC-6.12)}; las coordenadas de $P$ y $Q$ en \textup{(GMC-6.9)} dan
\textup{(GMC-6.13)}--\textup{(GMC-6.14)}. El argumento es reversible.
\end{proof}

\begin{bridgebox}
La clasificación \textup{(GMC-6.9)} proporciona la realización operatorial del puente. Convierte la
incompatibilidad suma--producto en un objeto operatorial completamente
calculable: cuatro sectores clásicos y dos bloques matriciales irreducibles.
El valor $1/3$ reaparece aquí como ángulo principal operatorial, no sólo como
solapamiento MUB; $5/7$ es el segundo valor genérico de $PQP$ en el sector de
suma--producto. Ninguno debe identificarse con la contracción $5/9$ del bloque
central $B_c$. Son invariantes de objetos distintos, coordinados sin
identificarlos por semejanza decimal.
\end{bridgebox}

\subsection{Separación espectral \texorpdfstring{$4$}{4} e intervalo matricial libre}

Para el bloque de \textup{(GMC-2.4)},
$B_c=7I+2R$ con $R=R^*$ no escalar, $R^2=I$ y
$\operatorname{Spec}(R)=\{-1,+1\}$,
$C^*(B_c)=C^*(R)\cong\mathbb C^2$. Su rango matricial en nivel $s$ es
\[
 \mathcal W_s(B_c)=
 \{\Phi(B_c):\Phi:C^*(B_c)\to M_s\text{ UCP}\}.
\]

\begin{theorem}[Intervalo libre del bloque central]
\label{p12-83-c22-s1:gmc:thm:intervalo-bc}
Para todo $s\ge1$,
\[
 \boxed{\ 
 \mathcal W_s(B_c)
 =\{Y=Y^*:5I_s\preceq Y\preceq9I_s\}.
 \ }
 \tag{GMC-6.15}
\]
\end{theorem}

\begin{proof}
Un mapa UCP envía la simetría $R$ a una contracción autoadjunta
$-I_s\preceq X\preceq I_s$ y, por tanto,
$\Phi(B_c)=7I_s+2X$ pertenece al intervalo. Recíprocamente, dado tal $X$,
los efectos $(I_s\pm X)/2$ definen un mapa UCP desde $\mathbb C^2$ que envía
$R$ a $X$.
\end{proof}

Así, $7$ es el centro, $2$ el radio operatorial y $4$ el diámetro espectral
del intervalo libre. El borde positivo, el borde del intervalo matricial y la
hipérbola de determinante uno de \textup{(GMC-2.6)} son tres nociones distintas que
ahora pueden coordinarse sin confundirse.

La doble proyección no reside en esta clasificación temprana. Su formulación
desde un antecedente espectral común se publica, con sus hipótesis y límites,
en la sección dedicada del capítulo rector de doble proyección.
\fi

\ifdefined\HMTGMCDoubleProjection
\subsection{Antecedente común de lectores: refinamiento espectral y límite de
dilatación conjunta}

Sea $X_m$ un conjunto finito de historias y sean
\[
 R_m:X_m\to\mathcal A_m,
 \qquad
 Q_m:X_m\to\mathcal E_m
 \tag{GMC-6.16}
\]
dos lectores: por ejemplo, uno arquimediano y otro excepcional. Cada uno
define una \PVM en $C(X_m)$ mediante las funciones características de sus
fibras.

\begin{proposition}[Antecedente espectral común]
\label{p12-83-c22-s1:gmc:prop:doble-proyeccion-conmutativa}
Las dos \PVM de \textup{(GMC-6.16)} conmutan y poseen el refinamiento común
\[
 P_{a,e}=1_{\{R_m=a\}}1_{\{Q_m=e\}}
 =1_{\{x:R_m(x)=a,\ Q_m(x)=e\}}.
 \tag{GMC-6.17}
\]
Si los lectores son compatibles con las truncaciones, los refinamientos son
compatibles en $m$ y determinan dos lectores sobre $C(X_\infty)$ con una
\PVM cilíndrica común.
\end{proposition}

\begin{proof}
Todas las funciones pertenecen a la misma álgebra conmutativa. El producto
de los indicadores es el indicador de la intersección. Sumar en $e$ recupera
la \PVM de $R_m$ y sumar en $a$ recupera la de $Q_m$. La compatibilidad pasa
al límite por el \Cref{p12-83-c22-s4:gmc:thm:dilatacion-diagonal}.
\end{proof}

Este resultado formaliza el antecedente común de la doble proyección. No
demuestra que el par de lectores sea conjuntamente inyectivo: para ello debe
separar puntos de $X_\infty$. Tampoco construye automáticamente un antecedente
no conmutativo común.

\begin{proposition}[Incompatibilidad de dos mediciones proyectivas mutuamente insesgadas]
Dos \PVM qutrit mutuamente insesgadas no admiten una medición conjunta aguda
en $M_3(\mathbb C)$.
\end{proposition}

\begin{proof}
Dos \PVM finitas son conjuntamente medibles por una \PVM común si y sólo si
conmutan. Proyectores de contextos MUB distintos tienen solapamiento $1/3$ y
no conmutan.
\end{proof}

\enlargethispage{3\baselineskip}
La doble proyección HMT no debe identificarse, por tanto, con dos contextos
MUB agudos sobre el mismo qutrit. Puede realizarse en la diagonal común de
historias, mediante efectos ruidosos, o en un antecedente no conmutativo más
amplio con dos canales UCP. La última posibilidad exige inclusión de rangos
matriciales, positividad de Choi y equivarianza respecto de la dinámica TPK.

\begin{workbox}[fontupper=\footnotesize,fonttitle=\bfseries\footnotesize,
  boxsep=1pt,left=4pt,right=4pt,top=0pt,bottom=0pt,
  toptitle=1pt,bottomtitle=1pt]
El antecedente no conmutativo se construye en
\Cref{sec:cpe-antecedente-ucp}: el sistema $\mathcal S_\infty$ y los mapas UCP
$\Phi_{\rm arq}$ y $\Phi_{\rm exc}$ reproducen simultáneamente las dos
proyecciones HMT y entrelazan la dinámica TPK en el dominio equivariante
declarado. La \Cref{p12-83-c22-s1:gmc:prop:doble-proyeccion-conmutativa}
proporciona su antecedente diagonal.
\end{workbox}
\fi


\section{Convexidad matricial genealógica: objetos, dominios y diccionario
tipado}

\subsection{Subestructuras, objeto global y proyecciones}

La pregunta que organiza el trabajo de De las Cuevas--Netzer es cómo una
colección de mediciones locales puede componer un objeto bidimensional y hasta
qué punto ese objeto es la compresión de una totalidad ideal. Una medición
aislada ---una \POVM--- admite dilatación de Naimark o Stinespring. Sin
embargo, una cuadrícula contiene simultáneamente una \POVM en cada fila y en
cada columna. Pedir una dilatación común que haga proyectivas todas esas
mediciones es una condición global mucho más fuerte. Hay cuadrados mágicos
cuánticos que no la satisfacen, incluso en el nivel matricial no conmutativo
más pequeño relevante \citep{delascuevas2020qms}.

HMT formula otro problema de partes y totalidades. La coordenada visible de un
estado no conserva necesariamente la ruta que la produjo. APP distingue las
evaluaciones suma y producto sobre un soporte común; TRIT conserva orientación
y acarreo; TPK conserva composición de caminos, retorno, holonomía y memoria.
El continuo surge de prolongar coherentemente las historias finitas, y la
doble proyección publica dos sombras diferentes de un mismo estado
enriquecido. Las dos teorías preguntan, por tanto, qué se pierde al comprimir.
La diferencia es que una organiza la pérdida por nivel matricial y la otra por
profundidad genealógica.

\begin{intuitionbox}
De las Cuevas pregunta: «¿puedo ver esta cuadrícula imperfecta como la
compresión de una cuadrícula proyectiva mayor?». HMT pregunta: «¿qué historia
completa fue comprimida para producir esta salida?». El puente responde que
ambas preguntas pueden formularse sobre la misma familia, con dos índices
independientes: profundidad $m$ y dimensión de observación $s$.
\end{intuitionbox}

\subsection{Cuadrados cuánticos, sistemas de operadores y convexidad matricial}

Sean $n,s\ge1$. Un cuadrado mágico cuántico de tamaño exterior $n$ e
interior $s$ es una matriz
\[
 A=(A_{ij})_{i,j=0}^{n-1}\in\Mat_n(\Her_s(\mathbb C))
\]
tal que
\[
 A_{ij}\succeq0,\qquad
 \sum_{j=0}^{n-1}A_{ij}=I_s,\qquad
 \sum_{i=0}^{n-1}A_{ij}=I_s.
 \tag{GMC-1.1}
\]
Cada fila y cada columna es una \POVM. Si todas las entradas son proyecciones,
$A$ es una matriz de permutación cuántica; si, además, esas proyecciones
conmutan, $A$ es una matriz de permutación cuántica conmutativa
\citep{delascuevas2020qms,delascuevas2023latin}.

Un \QMS es \emph{semiclásico} si existe una \POVM
\(\{Q_\sigma\}_{\sigma\in S_n}\) tal que
\[
 A=\sum_{\sigma\in S_n}P_\sigma\otimes Q_\sigma,
 \tag{GMC-1.2}
\]
donde $P_\sigma$ es la matriz de permutación clásica asociada a
$\sigma$. La fórmula no exige que los efectos $Q_\sigma$ conmuten. Dice
que toda la cuadrícula se obtiene reordenando una misma medición primaria.

Un cuadrado latino cuántico de orden $n$ es una cuadrícula de vectores
unitarios $v_{ij}\in\mathbb C^n$ cuyas filas y columnas son bases
ortonormales. Los proyectores de rango uno
\[
 P_{ij}=|v_{ij}\rangle\langle v_{ij}|
\]
forman un \QMS proyectivo. Los \QLS de clase fácil son los producidos por un cuadrado
latino clásico y una base ortonormal fija. De las Cuevas, Netzer y
Valentiner-Branth demuestran que los \QLS semiclásicos son exactamente los de
clase fácil y que la envolvente matricial convexa de los \QLS contiene más que los
cuadrados semiclásicos \citep{delascuevas2023latin}.

La operación matricial convexa elemental es
\[
 X=\sum_{r=1}^{k}V_r^*X^{(r)}V_r,
 \qquad
 \sum_{r=1}^{k}V_r^*V_r=I_s.
 \tag{GMC-1.3}
\]
Permite combinar objetos de tamaños internos diferentes. En lenguaje
operatorial, las compresiones y sumas directas se organizan mediante
aplicaciones unitales completamente positivas. El conjunto de todos los
valores matriciales de un sistema de operadores es su rango matricial.

\begin{importbox}
Los resultados rectores que se usan son: un \QMS se dilata a una matriz de
permutación cuántica con entradas conmutativas si y sólo si es semiclásico;
para $n\ge3$ y $s\ge2$ existen \QMS no semiclásicos, y para todo $n\ge3$ existen
\QMS sobre $M_2(\mathbb C)$ que no se dilatan a ninguna matriz de
permutación cuántica \citep{delascuevas2020qms}. La dilatabilidad simultánea
es, por tanto, una propiedad sustantiva, no una consecuencia automática de
que cada fila y columna se dilate por separado.
\end{importbox}

\subsection{Correspondencia tipada entre el estado genealógico HMT–MD y los sistemas de operadores: profundidad, truncamiento, monodromía y lectura}

En HMT una salida visible no agota el estado. Una sección finita
porta, según el dominio, datos del tipo
\[
 x_m=(\text{prefijo},\text{hoja},\text{orientación},\text{acarreo},
       \text{frontera},\text{ruta},\text{memoria},\text{calibre}).
 \tag{GMC-1.4}
\]
La profundidad $m$ no es una dimensión de Hilbert. Cuenta cuánto de la
genealogía se ha resuelto. Dos estados pueden compartir coordenada visible y
diferir en $m$, en hoja o en memoria.

APP proporciona el ortograma nonádico y las hojas suma y producto. TRIT
registra la orientación ternaria y el acarreo local. TPK es el núcleo
topológico de fase y la categoría de rutas que compone esos datos. En una
realización finita, las rutas actúan por operadores sobre un espacio de
Hilbert, pero la representación no sustituye la genealogía: la transporta.

El lenguaje de HMT contiene además tres operaciones que no están presentes en
la definición desnuda de un \QMS:

\begin{enumerate}
  \item \emph{truncación}: olvidar la continuación posterior a profundidad
  $m$, conservando el prefijo y agregando sus prolongaciones;
  \item \emph{monodromía}: retornar a la misma fase visible con memoria no
  trivial;
  \item \emph{lectura}: publicar una coordenada matemática o física sin
  identificarla con la sección completa.
\end{enumerate}

\Needspace{.74\textheight}
\subsection{Correspondencias directas e inversas entre los formalismos}

\begin{longtable}{@{}L{.23\textwidth}L{.29\textwidth}L{.37\textwidth}@{}}
\toprule
Lenguaje matricial & Lenguaje HMT & Correspondencia tipada\\
\midrule
\endfirsthead
\toprule
Lenguaje matricial & Lenguaje HMT & Correspondencia tipada\\
\midrule
\endhead
Tamaño exterior $n$ & cuadrícula publicada & número de filas y columnas del
objeto mágico; no es profundidad ni dimensión interna\\
Tamaño interior $s$ & carta operatorial & dimensión de los efectos o de la
compresión matricial\\
Profundidad $m$ & genealogía & índice nuevo que no se reduce a $n$ ni a
$s$\\
\POVM & familia de lectores positivos & efectos que suman la unidad; la suma
expresa normalización, no memoria completa\\
\PVM & contexto espectral & familia de proyecciones ortogonales que resuelve
la identidad\\
\UCP & lector matricial & mapa que conserva unidad y positividad a todos los
niveles\\
Compresión $V^*(\cdot)V$ & sombra & paso de una realización mayor a una
carta observacional menor\\
Dilatación & antecedente operatorio & representación multiplicativa común
antes de comprimir\\
Semiclasicidad & una \POVM primaria común reordenada & propiedad formal; no
equivale a conmutatividad de todas las celdas\\
No semiclasicidad & composición genuinamente bidimensional & no puede
reducirse al reordenamiento de una \POVM primaria común\\
Sistema de operadores & espacio de lectores & subespacio autoadjunto con la
unidad que organiza todas las cartas UCP\\
Rango matricial & conjunto de sombras a todos los $s$ & familia libre de
evaluaciones internas\\
Límite inductivo & cierre de refinamientos operatorios & conservación de la
unidad a través de toda la filtración nonádica\\
Límite inverso & sección infinita compatible & historia que conserva todos
sus prefijos\\
\bottomrule
\end{longtable}

\subsection{Extensión de la bigraduación a 3 índices}

El formalismo de cuadrados mágicos separa $(n,s)$. HMT añade $m$. El
objeto total debe escribirse
\[
 \mathfrak M^{(n)}_{m,s},
 \tag{GMC-1.5}
\]
y no simplemente \(\mathfrak M_n\) o \(\mathfrak M_s\). Cada índice responde
a una pregunta diferente:

\begin{center}
\begin{tabularx}{.92\textwidth}{@{}c L{.25\textwidth}Y@{}}
\toprule
Índice & Pregunta & Operación asociada\\
\midrule
$m$ & ¿cuánta historia se conserva? & refinamiento y truncación\\
$n$ & ¿cómo se compone la cuadrícula? & filas, columnas y permutaciones\\
$s$ & ¿desde qué nivel matricial se observa? & suma directa y compresión UCP\\
\bottomrule
\end{tabularx}
\end{center}

Para APP se fija $n=9$. Entonces la familia
$\{\mathfrak M^{(9)}_{m,s}\}_{m,s}$ está bigraduada por $(m,s)$. Si se
comparan APP, los cuatro \QLS de orden tres y el atlas de doce direcciones, el
índice $n$ vuelve a variar y la estructura correcta es trigraduada.

\begin{controlbox}
Llamar «bigraduación $(m,s)$» no borra $n$: significa que se estudia una
familia con tamaño exterior fijado. Llamar «$9\times9$» a APP tampoco fija
el tamaño interior de sus operadores. Las matrices internas $6\times6$, la
carta qutrit $s=3$ y la cuadrícula exterior $n=9$ son objetos de tipos
distintos.
\end{controlbox}

\subsection{Bidireccionalidad del intercambio matemático}

La dirección De las Cuevas--Netzer $\longrightarrow$ HMT aporta criterios que HMT no
tenía formulados con este grado de separación: positividad celda a celda,
normalización de todas las filas y columnas, semiclasicidad, frontera de la
envolvente matricial convexa, dilatación simultánea y extremos de Arveson. El
resultado es un conjunto de pruebas y falsadores para cualquier alegación de
«cuadrado cuántico».

La dirección HMT $\longrightarrow$ De las Cuevas--Netzer añade estructura al objeto que
se comprime: profundidad, ruta, no escisión del acarreo, monodromía, dinámica,
calibre, doble proyección y lectores metrológicos. El problema deja de ser
sólo si existe una dilatación y pasa a preguntar si existe una dilatación
\emph{genealógicamente fiel}: una que no identifique dos rutas distintas sólo
porque producen el mismo efecto visible.

\begin{bridgebox}
La noción nueva que organiza el artículo es la \emph{dilatación simultánea
genealógica}: una representación multiplicativa común que dilata no sólo varias
filas y columnas a profundidad fija, sino todas las profundidades compatibles,
y que especifica qué parte de la memoria queda en la representación y qué
parte se pierde en la compresión. El teorema exacto se demuestra en
\Cref{p12-83-c22-s4:gmc:chap:bigraduacion}.
\end{bridgebox}


\section[Sistemas operatoriales de incidencia]{Sistemas operatoriales de
incidencia: cuadrados latinos cuánticos, contextos qutrit y dilatación
genealógica común}
\label{p12-83-c22-s3:gmc:chap:qms}

El objeto exterior de APP tiene orden nueve, pero el número nueve no basta
para convertirlo en un cuadrado mágico cuántico. Hay que exhibir 81 efectos
positivos y verificar 18 normalizaciones operatoriales. Este capítulo realiza
esa tarea para dos sombras de APP y para el atlas completo de doce
direcciones. También determina su clase exacta en la taxonomía de De las
Cuevas--Netzer.

\subsection{4 cuadrados latinos cuánticos cíclicos de clase operatorial fácil}

Sean
\[
 \mathcal P_a=\{P_{a,0},P_{a,1},P_{a,2}\},
 \qquad a\in\mathbb P^1(\mathbb F_3),
 \tag{GMC-4.1}
\]
las cuatro \PVM qutrit del sistema de Weyl. Para cada $a$, elija vectores
unitarios $v_{a,k}$ con
$P_{a,k}=|v_{a,k}\rangle\langle v_{a,k}|$ y defina
\[
 L^{(a)}_{ij}=v_{a,i+j\bmod3}.
 \tag{GMC-4.2}
\]
Cada fila y cada columna recorre una vez la base
$(v_{a,0},v_{a,1},v_{a,2})$. Por tanto, $L^{(a)}$ es un \QLS de orden tres.
Como se obtiene de la tabla latina cíclica y de una base ortonormal fijada, es de clase fácil y su
\QMS proyectivo es semiclásico
\citep{delascuevas2023latin}.

\begin{proposition}[Obstrucción MUB al apilamiento latino]
\label{p12-83-c22-s3:gmc:prop:no-go-mub}
No se pueden tomar dos filas de contextos distintos
$\mathcal P_a$ y $\mathcal P_b$, $a\ne b$, como filas de un mismo \QLS de
orden tres conservando sus columnas alineadas.
\end{proposition}

\begin{proof}
La condición de columna de un \QLS exigiría ortogonalidad entre los dos
vectores de una columna. Pero el insesgamiento mutuo da
\[
 |\langle v_{a,k},v_{b,\ell}\rangle|^2
 =\Tr(P_{a,k}P_{b,\ell})=\frac13
\]
para todo $k,\ell$. El solapamiento no es cero.
\end{proof}

La obstrucción no niega la compatibilidad cuántica de las cuatro bases.
Determina que la estructura operatorial adecuada para usarlas simultáneamente no es un
\QLS de rango uno: hay que pasar a efectos ponderados y a un \QMS.

\subsection{Lema cíclico de cuadrados mágicos cuánticos asociados a una medida operatorial positiva}

\begin{lemma}[Cuadrado cíclico asociado a una \POVM]
\label{p12-83-c22-s3:gmc:lem:qms-ciclico}
Sea $E=(E_0,\ldots,E_{n-1})$ una \POVM sobre $\mathbb C^s$. Entonces
\[
 A_{ij}=E_{i+j\bmod n},
 \qquad 0\le i,j<n,
 \tag{GMC-4.3}
\]
es un \QMS de tamaños $(n,s)$. Además es semiclásico y admite una
dilatación simultánea a un \QPM conmutativo.
\end{lemma}

\begin{proof}
Cada entrada es positiva. Toda fila y toda columna es una permutación de los
$n$ efectos, de modo que su suma es $I_s$. Para $t\in\mathbb Z/n\mathbb Z$,
sea $\sigma_t(i)=t-i$ y $P_{\sigma_t}$ su matriz de permutación. Entonces
\[
 A=\sum_{t=0}^{n-1}P_{\sigma_t}\otimes E_t,
 \tag{GMC-4.4}
\]
que es una descomposición semiclásica explícita. Si
$E_t=V^*Q_tV$ es una dilatación de Naimark de la \POVM a una \PVM
$(Q_t)_t$, las entradas
$\widehat A_{ij}=Q_{i+j}$ forman un \QPM conmutativo y
$A_{ij}=V^*\widehat A_{ij}V$ simultáneamente para las $n^2$ celdas.
\end{proof}

Esta construcción pertenece a la clase de familias latinas asociadas a POVM
(\emph{POVM Latin-square friends})
de De las Cuevas y Netzer \citep{delascuevas2023latin}. La novedad del puente no es el lema
abstracto, sino las \POVM que HMT proporciona, su calibración nonádica y la
separación exacta entre lo que cada cuadrado conserva y lo que todavía
comprime.

\subsection{Representación aditiva de APP de tamaños \((9,3)\) y su codominio de incidencia}

Fijemos uno de los cuatro contextos, por ejemplo
$\mathcal P_0=\{P_{0,0},P_{0,1},P_{0,2}\}$. Para $t\in\mathbb Z/9\mathbb Z$
definimos
\[
 E_t^{\Sigma}=\frac13P_{0,t\bmod3},
 \qquad
 A_{ij}^{\Sigma}=E_{i+j\bmod9}^{\Sigma}.
 \tag{GMC-4.5}
\]
Cada proyección aparece tres veces, luego
\[
 \sum_{t=0}^{8}E_t^{\Sigma}
 =\frac13\sum_{k=0}^2 3P_{0,k}=I_3.
\]
Por el \Cref{p12-83-c22-s3:gmc:lem:qms-ciclico}, $A^\Sigma$ es un \QMS semiclásico de tamaños
$(9,3)$. Es una transducción directa de la ley latina
$\Sigma(i,j)=i+j\bmod9$ a una lectura qutrit. Conserva el tamaño exterior y
la hoja aditiva; deliberadamente no distingue las tres elevaciones Hensel de
una misma clase ternaria.

\subsection{\(3\) contextos en un cuadrado $(9,3)$}

Elijamos tres de las cuatro \PVM qutrit, indexadas por $a=0,1,2$, y ordenemos
sus proyectores mediante $t=3a+k$:
\[
 E_{3a+k}^{(9)}=\frac13P_{a,k},
 \qquad a,k\in\mathbb Z/3\mathbb Z.
 \tag{GMC-4.6}
\]
Como
\[
 \sum_{a,k}E_{3a+k}^{(9)}
 =\frac13\sum_{a=0}^2\sum_{k=0}^2P_{a,k}=I_3,
\]
la fórmula
\[
 A_{ij}^{(9)}=E_{i+j\bmod9}^{(9)}
 \tag{GMC-4.7}
\]
produce otro \QMS de tamaños $(9,3)$.

\begin{theorem}[Cuadrado qutrit de orden exterior nueve]
\label{p12-83-c22-s3:gmc:thm:qms-9-3}
El cuadrado $A^{(9)}$ de \textup{(GMC-4.7)} satisface:

\begin{enumerate}
 \item sus 81 celdas son efectos positivos de rango uno y traza $1/3$;
 \item todas las filas y columnas suman $I_3$;
 \item contiene entradas que no conmutan;
 \item es semiclásico mediante
 \[
 A^{(9)}=\sum_{t=0}^{8}P_{\sigma_t}\otimes E_t^{(9)};
 \tag{GMC-4.8}
 \]
 \item admite una dilatación simultánea conmutativa de sus 81
 entradas.
\end{enumerate}
\end{theorem}

\begin{proof}
Los dos primeros puntos y la descomposición siguen del
\Cref{p12-83-c22-s3:gmc:lem:qms-ciclico}. Si $a\ne b$, dos proyectores MUB de rango uno no
conmutan salvo coincidencia o ortogonalidad, imposibles porque su solapamiento
es $1/3$; por tanto algunos
$[E_{3a+k}^{(9)},E_{3b+\ell}^{(9)}]$ son no nulos. La última afirmación es la
dilatación construida en la prueba del lema.
\end{proof}

\begin{bridgebox}
Este teorema cierra una existencia que el puente preliminar había dejado
abierta: ya hay 81 efectos explícitos con $n=9$, $s=3$ y normalizaciones
exactas. Al mismo tiempo demuestra una separación conceptual fértil:
\[
 \text{celdas no conmutativas}
 \quad\centernot\Longrightarrow\quad
 \text{no semiclasicidad}.
\]
Aquí las celdas no conmutan, pero toda la cuadrícula procede de reordenar una
sola \POVM primaria.
\end{bridgebox}

\subsection{El atlas completo en un cuadrado \texorpdfstring{$(12,3)$}{(12,3)}}

La línea proyectiva sobre el anillo nonádico tiene doce puntos:
\[
 |\mathbb P^1(\mathbb Z/9\mathbb Z)|=12.
 \tag{GMC-4.9}
\]
La reducción módulo tres induce
\[
 r:\mathbb P^1(\mathbb Z/9\mathbb Z)
 \longrightarrow\mathbb P^1(\mathbb F_3),
 \tag{GMC-4.10}
\]
con cuatro fibras de cardinal tres. Fijada una calibración
\[
 \Xi:\mathbb P^1(\mathbb Z/9\mathbb Z)
 \xrightarrow{\sim}
 \{(a,k):a\in\mathbb P^1(\mathbb F_3),\ k\in\mathbb F_3\}
 \tag{GMC-4.11}
\]
que respeta la fibra de $r$, se asocia a cada dirección $\ell$ el proyector
$P_{\Xi(\ell)}$. Al ordenar las doce direcciones como $t=0,\ldots,11$, sea
\[
 F_t=\frac14P_{\Xi(\ell_t)},
 \qquad
 B_{ij}=F_{i+j\bmod12}.
 \tag{GMC-4.12}
\]

\begin{theorem}[Cuadrado qutrit del atlas nonádico completo]
\label{p12-83-c22-s3:gmc:thm:qms-12-3}
La matriz $B$ es un \QMS de tamaños $(12,3)$, contiene los doce proyectores
calibrados de los cuatro contextos qutrit y es semiclásica:
\[
 B=\sum_{t=0}^{11}P_{\sigma_t}\otimes F_t.
 \tag{GMC-4.13}
\]
Admite una dilatación simultánea a un \QPM conmutativo.
\end{theorem}

\begin{proof}
Cada uno de los cuatro contextos suma $I_3$, así que
$\sum_tF_t=(1/4)\cdot4I_3=I_3$. Se aplica el
\Cref{p12-83-c22-s3:gmc:lem:qms-ciclico}.
\end{proof}

El factor $1/4$ no es decorativo. El atlas contiene cuatro resoluciones de la
identidad y, sin esa normalización, la suma sería $4I_3$. El atlas completo
es, por tanto, una medición de doce resultados, no una quinta base
ortonormal.

\subsection{Teoremas operatoriales y criterio de fidelidad genealógica APP--TPK}

Las tres construcciones anteriores tienen un estatuto exacto:

\begin{center}
\begin{tabularx}{.96\textwidth}{@{}L{.20\textwidth}L{.21\textwidth}Y@{}}
\toprule
Objeto & Teorema cerrado & Estructura HMT aún no transportada\\
\midrule
$A^\Sigma$, $(9,3)$ & hoja suma, positividad, normalización y
semiclasicidad & hoja producto, cociente y rutas\\
$A^{(9)}$, $(9,3)$ & tres contextos, no conmutatividad y dilatación común &
cuarto contexto, incidencia Hensel, gnomon y memoria TPK\\
$B$, $(12,3)$ & atlas completo $12\to4\times3$ y dilatación común &
equivarianza de la calibración, dinámica y monodromía\\
\bottomrule
\end{tabularx}
\end{center}

Ninguno de los dos últimos cuadrados es un \QLS o un \QPM: sus celdas son
proyectores escalados por $1/3$ o $1/4$, no proyecciones. Sí son cuadrados
mágicos cuánticos en la definición ortodoxa y sí poseen una dilatación
simultánea. Esta formulación resuelve la ambigüedad terminológica y conserva
la diferencia matemática entre efecto y proyección.

\begin{workbox}
El siguiente objetivo no es volver a probar que existe un \QMS de orden
nueve. Es construir uno \emph{APP--TPK fiel}: además de \textup{(GMC-4.3)}, debe
portar simultáneamente $\Sigma$, $\Pi$, cociente, orientación y ruta, y debe
satisfacer una ecuación de equivarianza con la acción de TPK. Un candidato
que no pase esa ecuación seguirá siendo una sombra válida, pero no el objeto
total.
\end{workbox}


\section[Bigraduación genealógico--matricial]{Bigraduación
genealógico--matricial a tamaño exterior fijo y trigraduación
\texorpdfstring{$(m,n,s)$}{(m,n,s)}: truncación, compresión y dilatación
simultánea}
\label{p12-83-c22-s4:gmc:chap:bigraduacion}

La separación de De las Cuevas--Netzer entre tamaño exterior $n$ y nivel
matricial $s$ resuelve una ambigüedad del lenguaje de cuadrículas. HMT exige
una segunda separación: la profundidad $m$ no es un tamaño interior. Este
capítulo demuestra que las direcciones $m$ y $s$ son compatibles y que una
familia coherente en toda profundidad posee una dilatación común.

\subsection{Sistema inverso de historias y sistema directo de observables}

Sea $(X_m,p_m)_{m\ge0}$ un sistema de conjuntos finitos no vacíos con mapas
sobreyectivos
\[
 p_m:X_{m+1}\longrightarrow X_m.
 \tag{GMC-3.1}
\]
En HMT, $X_m$ puede ser el conjunto completo de prefijos nonádicos de
longitud $m$ o un conjunto de supervivientes siempre que las restricciones
sean sobreyectivas en el dominio declarado. El límite inverso
\[
 X_\infty=\varprojlim_m X_m
\]
es el espacio compacto de historias compatibles. Cada $x\in X_m$ determina
el cilindro clopen
\[
 [x]_m=\{\xi\in X_\infty:\xi_m=x\}.
 \tag{GMC-3.2}
\]

Por dualidad, las álgebras diagonales
\[
 D_m=C(X_m)
\]
forman un sistema inductivo mediante
\[
 \jmath_m:D_m\longrightarrow D_{m+1},
 \qquad \jmath_m(f)=f\circ p_m.
 \tag{GMC-3.3}
\]
Su límite uniforme es $C(X_\infty)$: las funciones cilíndricas forman una
subálgebra autoadjunta unital que separa puntos y son densas por
Stone--Weierstrass.

\subsection{La familia a \(2\) grados}

Para $s\ge1$, definimos el espacio de estados matriciales
\[
 \mathcal K_{m,s}=\UCP(D_m,M_s(\mathbb C)).
 \tag{GMC-3.4}
\]
Si $e_x\in D_m$ es la función característica de $x\in X_m$, toda
$\Phi\in\mathcal K_{m,s}$ equivale a una \POVM
\[
 E_x^{(m)}=\Phi(e_x)\succeq0,
 \qquad
 \sum_{x\in X_m}E_x^{(m)}=I_s.
 \tag{GMC-3.5}
\]
La restricción genealógica es
\[
 \rho_m^s:\mathcal K_{m+1,s}\longrightarrow\mathcal K_{m,s},
 \qquad
 \rho_m^s(\Phi)=\Phi\circ\jmath_m.
 \tag{GMC-3.6}
\]
En términos de efectos,
\[
 E_x^{(m)}=\sum_{y\in p_m^{-1}(x)}E_y^{(m+1)}.
 \tag{GMC-3.7}
\]
Es exactamente la operación «olvidar la continuación pero sumar todas sus
prolongaciones». No es un promedio ni una eliminación de masa: conserva la unidad.

En la dirección matricial, si $d\ge0$,
$\Phi_r\in\mathcal K_{d,s_r}$ y
$V_r:\mathbb C^s\to\mathbb C^{s_r}$ satisfacen
$\sum_rV_r^*V_r=I_s$, su combinación matricial convexa es
\[
 \operatorname{mc}_{\boldsymbol V}(\Phi_1,\ldots,\Phi_k)(f)
 =\sum_{r=1}^kV_r^*\Phi_r(f)V_r
 \in\mathcal K_{d,s}.
 \tag{GMC-3.8}
\]

\begin{theorem}[Conmutación de genealogía y convexidad matricial]
\label{p12-83-c22-s4:gmc:thm:conmutacion-ms}
Para toda familia $\Phi_r\in\mathcal K_{m+1,s_r}$ y toda compresión como en
\textup{(GMC-3.8)},
\[
 \rho_m^s\!\left(
   \operatorname{mc}_{\boldsymbol V}(\Phi_1,\ldots,\Phi_k)
 \right)
 =
 \operatorname{mc}_{\boldsymbol V}\!\left(
   \rho_m^{s_1}(\Phi_1),\ldots,\rho_m^{s_k}(\Phi_k)
 \right).
 \tag{GMC-3.9}
\]
\end{theorem}

\begin{proof}
Para $f\in D_m$, ambos miembros valen
\[
 \sum_rV_r^*\Phi_r(\jmath_m f)V_r.
\]
La igualdad es functorial, no numérica: vale simultáneamente para todos los
niveles matriciales y para cualquier sistema de truncaciones.
\end{proof}

\begin{bridgebox}
La ecuación \textup{(GMC-3.9)} es el núcleo de la bigraduación $(m,s)$. Refinar la
historia y después comprimir produce exactamente la misma sombra que
comprimir cada lector y después olvidar la continuación. La profundidad y el
nivel matricial son independientes, pero forman un cuadrado conmutativo.
\end{bridgebox}

\begin{figure}[H]
\centering
\begin{tikzpicture}[node distance=2.5cm and 3.7cm,>=Latex]
\node[draw,rounded corners,fill=hmtlightblue,minimum width=3.3cm,
 minimum height=1cm] (a) {$\mathcal K_{m+1,s'}$};
\node[draw,rounded corners,fill=hmtlightblue,right=of a,minimum width=3.3cm,
 minimum height=1cm] (b) {$\mathcal K_{m+1,s}$};
\node[draw,rounded corners,fill=hmtlightgold,below=of a,minimum width=3.3cm,
 minimum height=1cm] (c) {$\mathcal K_{m,s'}$};
\node[draw,rounded corners,fill=hmtlightgold,right=of c,minimum width=3.3cm,
 minimum height=1cm] (d) {$\mathcal K_{m,s}$};
\draw[->,thick] (a)--node[above]{compresión} (b);
\draw[->,thick] (a)--node[left]{truncación} (c);
\draw[->,thick] (b)--node[right]{truncación} (d);
\draw[->,thick] (c)--node[below]{compresión} (d);
\node[hmtviolet] at ($(a)!0.5!(d)$) {conmuta};
\end{tikzpicture}
\caption{Las dos direcciones independientes de la bigraduación. El símbolo
$s'$ puede representar una suma directa de varios niveles internos.}
\label{p12-83-c22-s4:gmc:fig:cuadrado-ms}
\end{figure}

\subsection{Dilatación uniforme compatible con todas las profundidades del sistema proyectivo}

Una familia $(\Phi_m)_{m\ge0}$ es compatible si
\[
 \Phi_m=\Phi_{m+1}\circ\jmath_m.
 \tag{GMC-3.10}
\]
Equivale a exigir la ley de agregación \textup{(GMC-3.7)} para todos los prefijos.

\begin{theorem}[Dilatación simultánea genealógica, corredor diagonal]
\label{p12-83-c22-s4:gmc:thm:dilatacion-diagonal}
Toda familia compatible
$\Phi_m\in\UCP(C(X_m),M_s)$ determina una única aplicación
\[
 \Phi_\infty:C(X_\infty)\longrightarrow M_s
 \tag{GMC-3.11}
\]
que coincide con $\Phi_m$ en cada álgebra cilíndrica. Existen un espacio de
Hilbert $\mathcal K$, una isometría $V:\mathbb C^s\to\mathcal K$, una
representación $\pi:C(X_\infty)\to B(\mathcal K)$ y su medida espectral
proyectiva asociada $P$ sobre $X_\infty$ tales que
\[
 \Phi_\infty(f)=V^*\left(\int_{X_\infty}f(\xi)\,dP(\xi)\right)V.
 \tag{GMC-3.12}
\]
En particular, para todo $m$ y $x\in X_m$,
\[
 E_x^{(m)}=V^*P([x]_m)V.
 \tag{GMC-3.13}
\]
Una \PVM común dilata simultáneamente todas las \POVM finitas de la
genealogía. Si la dilatación de Stinespring se elige mínima, el triple
$(\pi,\mathcal K,V)$ es único hasta equivalencia unitaria; una vez fijada
$\pi$, su medida espectral $P$ es única.
\end{theorem}

\begin{proof}
La compatibilidad define una aplicación unital positiva sobre la unión de las
álgebras cilíndricas. Como el dominio es conmutativo, positividad implica
positividad completa, y la aplicación es contractiva. Por densidad se
extiende de forma única a $C(X_\infty)$. El teorema de Stinespring produce
$\Phi_\infty(f)=V^*\pi(f)V$ y garantiza que su realización mínima es única
hasta equivalencia unitaria. Toda representación de una álgebra conmutativa
es la integral respecto de una única \PVM $P$ asociada a esa representación.
Aplicar la fórmula a la función
característica del cilindro $[x]_m$ da \textup{(GMC-3.13)}.
\end{proof}

\begin{corollary}[Realización ambiental fiel]
La representación de \Cref{p12-83-c22-s4:gmc:thm:dilatacion-diagonal} puede elegirse fiel sobre
$C(X_\infty)$ sin cambiar ninguna compresión finita.
\end{corollary}

\begin{proof}
Si $(\pi,\mathcal K,V)$ es una dilatación y
$\pi_f:C(X_\infty)\to B(\mathcal K_f)$ una representación fiel, se toma
$\pi\oplus\pi_f$ y la isometría $V\oplus0$. La compresión sigue siendo
$\Phi_\infty$ y la representación ambiental ya no identifica dos funciones
cilíndricas distintas.
\end{proof}

La última construcción tipa la frase «conservar la genealogía». La
dilatación mínima conserva exactamente lo que ve el lector; la dilatación
ambiental fiel conserva además todas las distinciones del álgebra antes de
comprimir. La fidelidad de una dinámica TPK, sin embargo, requiere también
entrelazar sus flechas y no se deduce sólo de la fidelidad de la
representación diagonal.

\subsection{Torre UHF no conmutativa y compatibilidad matricial}

La torre UHF nonádica, su traza y sus cierres factoriales se demuestran en
\cref{chap:integracion-vn64-operatorial}. Aquí se toman como estructura de
partida explícita para demostrar un resultado distinto: la prolongación
simultánea de una familia compatible de aplicaciones UCP y su dilatación
común.

Sea
\[
 \mathcal A_m=M_{9^m}(\mathbb C),
 \qquad \iota_m(a)=a\otimes I_9,
 \qquad
 \mathcal A_\infty=\varinjlim(\mathcal A_m,\iota_m).
 \tag{GMC-3.14}
\]

\begin{theorem}[Dilatación simultánea sobre el límite UHF]
\label{p12-83-c22-s4:gmc:thm:dilatacion-uhf}
Si $\Psi_m:\mathcal A_m\to M_s$ son aplicaciones UCP y
\[
 \Psi_m=\Psi_{m+1}\circ\iota_m,
 \tag{GMC-3.15}
\]
existe una única aplicación UCP
$\Psi_\infty:\mathcal A_\infty\to M_s$ que las prolonga. Existe una
representación $\pi:\mathcal A_\infty\to B(\mathcal K)$ y una isometría
$V:\mathbb C^s\to\mathcal K$ tales que
\[
 \Psi_m(a)=V^*\pi(\iota_{m,\infty}(a))V
 \quad\text{para todo }m.
 \tag{GMC-3.16}
\]
La representación puede hacerse fiel por suma directa con una representación
fiel del límite UHF.
\end{theorem}

\begin{proof}
La compatibilidad define $\Psi_\infty$ en el límite inductivo algebraico. Las
aplicaciones UCP son completamente contractivas, por lo que se extiende a la
completación $C^*$. Stinespring da \textup{(GMC-3.16)}; la suma directa demuestra la
última afirmación.
\end{proof}

Al restringir \Cref{p12-83-c22-s4:gmc:thm:dilatacion-uhf} a las diagonales canónicas se recupera
el caso nonádico completo $X_m=(\mathbb Z/9\mathbb Z)^m$ del
\Cref{p12-83-c22-s4:gmc:thm:dilatacion-diagonal}; esta restricción no cubre por sí sola un
sistema arbitrario de supervivientes. En el sistema UHF nonádico canónico de
\textup{(GMC-3.14)}, la representación GNS de su traza única tiene clausura débil
isomorfa a $\Rhyper$. Así, la bigraduación no es un apéndice combinatorio: une
los cilindros de historias, las cartas matriciales finitas y el continuo
operatorial de la realización de von Neumann.

\subsection{Trigraduación y niveles de simultaneidad}

Cuando aparece una cuadrícula, se incorpora el tamaño exterior $n$:
\[
 \mathfrak M_{m,s}^{(n)}.
 \tag{GMC-3.17}
\]
La estructura es trigraduada si $n$ varía y bigraduada por $(m,s)$ si se fija
$n$. «Simultáneo» puede significar tres exigencias crecientes:

\begin{enumerate}
 \item \emph{por filas y columnas}: una dilatación común de las $2n$ \POVM
 de un cuadrado mágico;
 \item \emph{por profundidades}: una dilatación común de todos los niveles
 $m$, demostrada en los Teoremas~\ref{p12-83-c22-s4:gmc:thm:dilatacion-diagonal}
 y~\ref{p12-83-c22-s4:gmc:thm:dilatacion-uhf};
 \item \emph{por genealogía dinámica}: además, un entrelazador que preserve
 ruta, monodromía y composición TPK.
\end{enumerate}

Las dos primeras poseen teoremas exactos en sus dominios. La tercera es la
noción más fuerte que HMT aporta al marco de De las Cuevas--Netzer. La
construcción del sistema de operadores común y su ecuación de equivarianza se
demuestran en \Cref{sec:cpe-antecedente-ucp}; la dilatación estática constituye
solamente uno de sus antecedentes.

\begin{controlbox}
El teorema no afirma que cualquier colección inconexa de efectos a distintas
profundidades tenga una dilatación común. La hipótesis decisiva es la
compatibilidad \textup{(GMC-3.10)} o \textup{(GMC-3.15)}. En HMT esa compatibilidad es la
forma operatorial de «cada prefijo es exactamente la restricción de sus
continuaciones».
\end{controlbox}


\section{Conos matriciales, rangos y extensiones completamente positivas}
\label{p12-83-c22-s5:gmc:chap:beyond}

El programa reciente de De les Coves, van der Eyden y Netzer amplía los
sistemas de operadores a sistemas cónicos definidos sobre categorías más
generales \citep{delescoves2026beyond}. El marco no convierte cualquier
categoría en teoría cuántica: exige elegir conos, dualidades, compresiones y
mapas completamente positivos adecuados. HMT aporta un banco especialmente
rico para realizar esa elección.

\subsection{Realizaciones en sistemas de operadores ordinarios}

Las familias siguientes están completamente dentro del formalismo estándar:

\begin{itemize}
 \item $\mathcal S_d=\Span\{I,P_{\Sigma,d},P_{\Pi,d}\}$ y sus rangos
 matriciales;
 \item las álgebras diagonales $C(X_m)$ y la torre UHF
 $M_{9^m}\hookrightarrow M_{9^{m+1}}$;
 \item los lectores UCP, \POVM, \PVM, canales CPTP y dilataciones de
 Stinespring;
 \item las familias bigraduadas $\mathcal K_{m,s}$.
\end{itemize}

Aquí no hace falta un formalismo categórico más amplio para afirmar
positividad, dilatación o convexidad matricial. De hecho, la clasificación
\textup{(GMC-6.9)} muestra que una parte sustancial de APP puede resolverse con
teoría operatorial finito-dimensional exacta.

\subsection{Extensión genealógica de los sistemas de operadores para rutas TPK}

Una ruta TPK no es sólo un operador entre dos espacios: porta origen, destino,
hoja, orientación, longitud, acarreo y una ley de composición. Si se olvidan
esas etiquetas y se conserva únicamente la matriz que representa la ruta,
pueden colapsar flechas distintas. Para incorporar el tipo completo se propone
una categoría $\mathsf C_{\rm HMT}$ cuyos objetos sean perfiles de registro y
cuyas flechas sean transportes admisibles. A cada tipo $c$ se asociaría un
cono de lectores positivos $A(c)$, y a cada transformación una acción
compatible sobre esos conos.

En el lenguaje de \emph{Beyond Operator Systems}, el objetivo es construir un
$S$-system: una extensión functorial del producto tensorial por el objeto base
a una familia de conos, con mapas CP definidos como transformaciones
naturales nivel a nivel. La teoría proporciona sistemas mínimos y máximos,
dualidad, teoremas de Choi--Kraus generalizados, representaciones como
intersecciones de $S$-hedros y una categoría $*$-autónoma
\citep{delescoves2026beyond}.

\begin{workbox}
El $S$-system intrínseco de TPK aún debe construirse. Faltan, como mínimo: el
funtor de conos, el objeto dualizante, las evaluaciones y coevaluaciones, y la
prueba de que la composición TPK respeta la positividad en todos los tipos.
El hecho de que TPK sea monoidal en su realización finita no demuestra por sí
solo que sea $*$-autónoma.
\end{workbox}

\subsection{Conos genealógicos HMT y ampliación de la teoría cónica}

Los ejemplos clásicos parten de un cono base y estudian sus sistemas mínimo y
máximo. HMT sugiere un cono base con memoria estratificada:
\[
 \mathcal C_m^{\rm HMT}
 =\operatorname{cone}\{\text{lectores positivos de historias a profundidad }m\}.
 \tag{GMC-8.1}
\]
Las truncaciones inducen mapas positivos
$\mathcal C_{m+1}^{\rm HMT}\to\mathcal C_m^{\rm HMT}$, mientras las
compresiones generan los niveles matriciales. El nuevo fenómeno a estudiar es
si los sistemas mínimo y máximo permanecen separados al tomar el límite en
$m$, y qué parte de esa separación mide pérdida de genealogía en vez de mera
no conmutatividad.

Esto propone tres invariantes para el programa general:

\begin{enumerate}
 \item \emph{profundidad de separación}: el menor $m$ en que los sistemas
 mínimo y máximo difieren;
 \item \emph{nivel matricial de detección}: el menor $s$ que exhibe esa
 diferencia;
 \item \emph{defecto de memoria}: la mínima información de ruta que hay que
 añadir para levantar un punto del sistema máximo al mínimo.
\end{enumerate}

Los dos primeros refinan la separación $(m,s)$; el tercero es genuinamente
HMT. Son definiciones propuestas y deben estudiarse primero en
$\mathfrak C_2^{\rm APP}$, donde \textup{(GMC-6.9)} permite cálculos exactos.

\subsection{Construcción explícita de cubos mágicos cuánticos y compatibilidad de sus márgenes}

El artículo de cuadrados latinos propone desarrollar cubos mágicos cuánticos.
El lema cíclico admite una extensión inmediata.

\begin{definition}[Cubo mágico cuántico]
Una familia
$C=(C_{ijk})_{0\le i,j,k<n}$ de efectos en $M_s$ es un cubo mágico cuántico
si toda línea paralela a cualquiera de los tres ejes suma $I_s$.
\end{definition}

\begin{theorem}[Cubo cíclico de una \POVM]
\label{p12-83-c22-s5:gmc:thm:cubo-ciclico}
Para toda \POVM $(E_t)_{t\in\mathbb Z/n\mathbb Z}$,
\[
 C_{ijk}=E_{i+j+k\bmod n}
 \tag{GMC-8.2}
\]
es un cubo mágico cuántico de tamaños exterior $n$ e interior $s$.
\end{theorem}

\begin{proof}
Fijadas dos coordenadas, al variar la tercera el índice
$i+j+k$ recorre una vez $\mathbb Z/n\mathbb Z$. Cada línea suma, por tanto,
$\sum_tE_t=I_s$. La positividad es heredada de la \POVM.
\end{proof}

Aplicando el teorema a $E^{(9)}$ de \textup{(GMC-4.6)} y a $F$ de
\textup{(GMC-4.12)} se obtienen cubos qutrit explícitos de órdenes nueve y doce. El
primero puede dedicar el tercer eje a una coordenada de fase; el segundo usa
el atlas nonádico completo. Para convertir ese eje en profundidad HMT real,
las capas deben además satisfacer
\[
 \rho_m(C^{m+1}_{ijk})=C^m_{ijk},
 \tag{GMC-8.3}
\]
y entrelazar TPK. La construcción cíclica cierra el cubo operatorial; la
fidelidad genealógica sigue siendo una condición adicional.

\subsection{Combinaciones mágicas de mapas CP}

Otra dirección abierta de la teoría consiste en trasladar las condiciones
mágicas a mapas completamente positivos mediante Choi--Jamiołkowski. Si
$J(\Phi_{ij})\succeq0$ son las matrices de Choi de una cuadrícula de mapas CP,
las condiciones
\[
 \sum_jJ(\Phi_{ij})=J(\id),
 \qquad
 \sum_iJ(\Phi_{ij})=J(\id)
 \tag{GMC-8.4}
\]
son una versión lineal de las normalizaciones mágicas. HMT aporta canales de
lectura, actualización y transporte; el problema es descomponerlos en una
cuadrícula que preserve a la vez normalización, hoja y ruta. El formalismo
cónico generalizado puede tratar codominios que no sean simplemente conos
PSD matriciales.

\subsection{Criterio de universalidad en el dominio genealógico}

El marco de Gonda, Reinhart, Stengele y De les Coves define un simulador en
una categoría \emph{objetivo--contexto} mediante
\[
 s:P\otimes C\longrightarrow T\otimes C,
 \tag{GMC-8.5}
\]
con un compilador funcional $s_T:P\to T$ y una reducción de contexto. La
universalidad es la existencia de una reducción laxa al simulador trivial,
no la mera riqueza o recurrencia del sistema
\citep{gonda2024universality}.

Una instancia HMT puede tomar:

\begin{itemize}
 \item $P$: prefijos, programas de ruta o semillas;
 \item $T$: observables matemáticos y físico-metrológicos que se pretende
 realizar;
 \item $C$: calibre, hoja, condiciones de frontera y aparato de lectura;
\item $s_T$: transformación de una genealogía en una salida;
 \item $s_C$: transporte del contexto y de la memoria residual.
\end{itemize}

Para demostrar universalidad habría que construir esos datos y la reducción
laxa. Para refutar una propuesta basta un monótono del preorden contextual.
El teorema de imposibilidad del marco afirma, en particular, que si la imagen
funcional del compilador tiene supremum estrictamente menor que el espacio de
objetivos para algún monótono compatible, el simulador no es universal.

\begin{bridgebox}
HMT sugiere una versión más fuerte: \emph{universalidad genealógicamente
fiel}. Además de reproducir el comportamiento del objetivo, la reducción debe
levantar a un funtor entre categorías de rutas y preservar composición,
retorno y registro cronológico enriquecido. Dos simuladores conductualmente equivalentes pueden dejar
de ser equivalentes cuando se exige conservar la procedencia. Esta distinción
ofrece una instancia nueva y no trivial para el programa categórico.
\end{bridgebox}

La unificación metrológica demostrada por HMT y la universalidad categórica
son afirmaciones distintas. La primera deriva lectores y relaciones dentro
del sistema; la segunda exige poder simular una clase declarada de objetivos y
contextos. Tiparlas por separado fortalece ambas: evita pedir a la metrología
una propiedad que no necesita y convierte la universalidad en un teorema
falsable.

\subsection{Contextualidad y no semiclasicidad}

Los cuadrados cíclicos de \Cref{p12-83-c22-s3:gmc:chap:qms} son semiclásicos. Por eso no
resuelven la dirección abierta sobre contextualidad o sobre QMS no
semiclásicos. APP, sin embargo, contiene el par no conmutativo $P_\Sigma,P_\Pi$
y el álgebra \textup{(GMC-6.9)}. La ruta correcta es construir efectos de celda que
dependan de ambos operadores y probar que no admiten la descomposición
\textup{(GMC-4.4)} con una \POVM sobre $S_n$.

La comprobación es un problema semidefinido: buscar
$Q_\sigma\succeq0$, $\sum_\sigma Q_\sigma=I$ y
\[
 A_{ij}=\sum_{\sigma:\,\sigma(i)=j}Q_\sigma.
 \tag{GMC-8.6}
\]
La factibilidad certifica semiclasicidad; la inviabilidad debe acompañarse de
un testigo dual. Una matriz no conmutativa o una gráfica visual no sustituyen
ese certificado.


\section{Estructura genealógica de APP--TRIT--TPK en nivel matricial}
\label{p12-83-c22-s6:gmc:chap:hmt-profundidad}

Este capítulo aísla las estructuras que intervienen en la transducción y
conserva explícitos sus tipos. El orden es
deliberado: soporte, diferencia suma--producto, memoria ternaria, categoría de
rutas, realización cuántica, cierre operatorial y lectura metrológica. Sólo
después se preguntará cuáles de esas estructuras forman un cuadrado mágico o
un sistema de operadores.

\subsection{Soporte común de APP y evaluación dual mediante las hojas suma y producto}

Sea
\[
 X_9=(\mathbb Z/9\mathbb Z)^2,
 \qquad
 \mathscr D=\{(1,0),(0,1),(-1,0),(0,-1)\}.
\]
El grafo de Cayley $\operatorname{Cay}(X_9,\mathscr D)$ tiene 81 vértices y
cuatro aristas cardinales orientadas desde cada vértice. La arquitectura APP
contiene todos los caminos finitos de ese grafo. La sigla se conserva como
nombre técnico y aquí no se usa para fijar una expansión histórica. A longitud
$r$ hay exactamente $81\cdot4^r$ caminos si se permiten retornos y
retrocesos. No hay una tabla aditiva y otra multiplicativa independientes:
hay un soporte único y dos evaluaciones del mismo punto,
\[
 \Sigma(i,j)=\operatorname{dr}(i+j),
 \qquad
 \Pi(i,j)=\operatorname{dr}(ij),
 \qquad
 \operatorname{dr}(a)=1+((a-1)\bmod9).
 \tag{GMC-2.1}
\]
La tabla de $\Sigma$ es un cuadrado latino clásico de orden nueve, porque
toda traslación permuta $\mathbb Z/9\mathbb Z$. La tabla de $\Pi$ no lo es:
sólo sus filas y columnas indexadas por las seis unidades
$\{1,2,4,5,7,8\}$ son permutaciones; $3,6,9$ pliegan clases.

\begin{resultbox}
APP ya contiene una estructura latina exacta, pero APP no se agota en ella.
La hoja suma es latina; la hoja producto registra el defecto de latinidad del
anillo local; los caminos conservan el orden de consulta, y la elevación
residuo--cociente conserva lo que la reducción modular ocultaría.
\end{resultbox}

Para cada celda se retienen residuo y cociente:
\[
 i+j=R^+_{ij}+9q^+_{ij},
 \qquad
 ij=R^\times_{ij}+9q^\times_{ij}.
 \tag{GMC-2.2}
\]
Los totales certificados son
\[
 \sum R^+=405,
 \quad \sum R^\times=459,
 \quad \sum q^+=45,
 \quad \sum q^\times=174,
\]
y por tanto
\[
 2025-810=1215=(459-405)+9(174-45)=54+9\cdot129.
 \tag{GMC-2.3}
\]
La diferencia visible $54$ y la diferencia de cociente $129$ no son dos
nombres para el mismo dato. Una vive en los residuos publicados; la otra, en
la memoria necesaria para recomponer el valor entero.

\subsection{Jerarquía espectral \(4\to2\to6\to54\) del bloque central de APP}

El bloque central de APP toma la forma
\[
 B_c=\begin{pmatrix}7&2\\2&7\end{pmatrix}
     =7I+2R
     =9P_++5P_-,
 \qquad
 R=\begin{pmatrix}0&1\\1&0\end{pmatrix},
 \quad P_\pm=\frac{I\pm R}{2}.
 \tag{GMC-2.4}
\]
De aquí se obtiene exactamente la jerarquía tipada inducida por la
descomposición espectral anterior:
\[
 \underbrace{9-5}_{4}
 \ ;\ 
 \underbrace{2}_{\text{acoplamiento fuera de la diagonal}}
 \ ;\ 
 \underbrace{51-45}_{6}
 \ ;\ 
 \underbrace{9\cdot6}_{54}.
 \tag{GMC-2.5}
\]
La primera cifra es el salto espectral primordial; la segunda, el acoplamiento
del bloque; la tercera, el exceso agregado del sector no unitario; la cuarta,
su transporte a la escala global de la tabla. No deben colapsarse por su
parentesco aritmético.

Como $\det B_c=45$, la normalización tiene determinante uno y pertenece al
grupo hiperbólico conectado:
\[
 \frac{B_c}{\sqrt{45}}
 =\exp(\eta R),
 \qquad
 \eta=\frac12\log\frac95,
 \qquad
 \frac59=e^{-2\eta}.
 \tag{GMC-2.6}
\]
La misma pareja espectral admite así una lectura contractiva y una lectura
lorentziana. Esta es la conexión precisa con convexidad e interior
hiperbólico que debe conservar el puente: la contracción $5/9$ es la sombra
de la separación de los dos idempotentes $P_\pm$, no una analogía verbal con
la curvatura.

\subsection{TRIT: estado visible, orientación, acarreo y memoria genealógica}

TRIT es la firma orientada $\{-1,0,+1\}$ obtenida de $\mathbb F_3$ mediante
$0\mapsto0$, $1\mapsto+1$, $2\mapsto-1$. Para tipar su elevación local,
definimos el alfabeto levantado
\[
 \mathcal L_1=\{\nu=r+3c:r,c\in\{-1,0,+1\}\}
 =\{-4,-3,\ldots,4\},
 \qquad
 \widehat\tau(\nu)=(r,c).
 \tag{GMC-2.7}
\]
La descomposición $\nu=r+3c$ es única en $\mathcal L_1$; $r$ es el residuo
visible y $c$ el acarreo. El símbolo $\nu$ es una amplitud levantada y no se
identifica, sin más datos, con la suma de dos trits visibles. En particular,
las amplitudes $3,0,-3$ publican el mismo residuo cero y dan tres antecedentes
distintos,
\[
 0^+=(0,+1),\qquad 0^0=(0,0),\qquad0^-=(0,-1).
 \tag{GMC-2.8}
\]
El estado completo puede contener
\[
 x=(t,\mu,o,h,c,\sigma,\lambda,g),
 \tag{GMC-2.9}
\]
donde $t$ es la firma visible, $\mu$ la hoja, $o$ la orientación, $h$ la
extensión, $c$ el acarreo, $\sigma$ la frontera, $\lambda$ la ruta y $g$ la
fase nonádica. La proyección $x\mapsto t$ elimina todas las demás
coordenadas. Éste es el punto estructural que HMT aporta a la teoría de
dilataciones: dos efectos observables iguales pueden poseer genealogías
distintas.

Las tres firmas admiten las álgebras reales
\[
 \mathbb A_\tau=\mathbb R[J]/(J^2+\bar\tau),
\]
de tipo elíptico ($J^2=-I$), parabólico ($J^2=0$) e hiperbólico
($J^2=+I$). En el último caso, $P_\pm=(I\pm J)/2$ separan dos direcciones
propias. El TRIT no es sólo el signo: es el observable mínimo de un transporte
que conserva etiqueta, orientación y memoria.

\subsection{TPK: categoría de rutas con memoria}

TPK, \emph{Topological Phase Kernel}, vive sobre la categoría libre de
caminos de APP y sobre su lectura bivariada suma--producto. Dos caminos que
terminan en la misma celda pueden diferir por palabra cardinal, enrollamiento,
hoja, cociente y holonomía. TPK conserva esa diferencia y compone rutas de
forma asociativa. El cociclo de acarreo
\[
 \kappa_9(a,b)=
 \frac{\operatorname{dr}(a)+\operatorname{dr}(b)
       -\operatorname{dr}(a+b)}9
\]
satisface
\[
 \kappa_9(a,b)+\kappa_9(a+b,c)
 =\kappa_9(b,c)+\kappa_9(a,b+c),
 \tag{GMC-2.10}
\]
y garantiza que la recuperación no depende de cómo se parentice una
concatenación.

El cociente fase--longitud posee objetos $(b,r)\in\mathbb Z/12\mathbb Z
\times\{0,\ldots,8\}$ y flechas
\[
 (b,r)\xrightarrow{\ell}
 \left(b+\left\lfloor\frac{r+\ell}{9}\right\rfloor,
       r+\ell\bmod9\right).
 \tag{GMC-2.11}
\]
Después de nueve iteraciones unitarias ($\ell=1$) retorna $r$ y avanza una fase
dodecafásica; después de $108=9\cdot12$ iteraciones unitarias retorna el cociente.
El estado enriquecido del TPK puede no
haber retornado, porque conserva memoria adicional. El cociente predictivo
mínimo de los observables publicados tiene 36 estados y fibras de cardinal
tres. Este hecho hace visible la diferencia entre igualdad de salida y
equivalencia de historia.

\begin{intuitionbox}
Un cuadrado mágico organiza simultáneamente filas y columnas. TPK organiza
simultáneamente posición, hoja, fase y composición. La traducción correcta no
consiste en llamar «cuántica» a toda tabla HMT, sino en construir un lector
positivo que preserve las identidades de APP y declarar, coordenada por
coordenada, qué memoria conserva y qué memoria comprime.
\end{intuitionbox}

\subsection{Realización hilbertiana y estructura cuántica}

La clausura cuántica finita de HMT construye, para cada profundidad
$d$, el espacio
\[
 \mathcal H_d=\ell^2((\mathbb Z/9\mathbb Z)^d),
 \qquad
 \mathcal A_d=B(\mathcal H_d)\cong M_{9^d}(\mathbb C).
 \tag{GMC-2.12}
\]
Para $u,v\in(\mathbb Z/9\mathbb Z)^d$, las traslaciones y fases nonádicas
actúan mediante
\[
 X_u|x\rangle=|x+u\rangle,
 \qquad
 Z_v|x\rangle=\omega^{v\cdot x}|x\rangle,
 \qquad \omega=e^{2\pi i/9},
\]
y satisfacen
\[
 Z_vX_u=\omega^{u\cdot v}X_uZ_v.
\]
La familia completa ---equivalentemente, los $d$ pares coordenados
$X_j,Z_j$--- genera $\mathcal A_d$ y actúa irreduciblemente. Para el carácter
central primitivo fijado por $\omega$, la representación finita de Heisenberg
queda determinada, salvo equivalencia unitaria. En esta realización se
construyen estados, observables autoadjuntos, medidas
espectrales, probabilidades de Born, dinámica unitaria y canales de Kraus; la
composición de registros se transporta monoidalmente al producto tensorial.
La representación hilbertiana finita no se presupone: se construye
explícitamente y transporta la genealogía del sistema.

En el cociente qutrit aparecen cuatro subálgebras abelianas maximales del
sistema de Weyl de dimensión tres. Sus cuatro bases espectrales son
mutuamente insesgadas:
\[
 \Tr(P_{a,k}P_{b,\ell})=
 \begin{cases}
  \delta_{k\ell},&a=b,\\[2pt]
  1/3,&a\ne b,
 \end{cases}
 \qquad a,b\in\mathbb P^1(\mathbb F_3).
 \tag{GMC-2.13}
\]
Los doce proyectores minimales $P_{a,k}$ son el atlas operatorial que más
adelante se calibra con las doce direcciones de
$\mathbb P^1(\mathbb Z/9\mathbb Z)$.

\subsection{De la torre nonádica al continuo operatorial}

La inclusión unital, la rama marcada, el límite UHF y los cierres de tipos
\(I_\infty\) y \(II_1\) se construyen en
\cref{chap:integracion-vn64-operatorial}. Sus fórmulas son premisas del puente
matricial y determinan la bigraduación.

La torre operatorial prolonga las cartas finitas mediante las
inclusiones unitarias y preservadoras de la traza normalizada
\[
 \iota_d:M_{9^d}(\mathbb C)\longrightarrow M_{9^{d+1}}(\mathbb C),
 \qquad a\longmapsto a\otimes I_9.
 \tag{GMC-2.14}
\]
En efecto, si $\tau_n=\Tr/n$, entonces
$\tau_{9^{d+1}}\circ\iota_d=\tau_{9^d}$; la traza ordinaria, en cambio, se
multiplica por nueve. 
Su límite inductivo es el álgebra UHF de tipo $9^\infty=3^\infty$. El cierre
en la representación GNS de la traza única es el factor hiperinfinito
$\Rhyper$ de tipo $\mathrm{II}_1$. Las inclusiones marcadas
$a\mapsto a\otimes p_j$ siguen una rama y conducen al corredor de compactos,
de tipo $\mathrm I_\infty$. El resultado distingue, por tanto, conservar
simultáneamente todas las ramas de seleccionar una historia concreta.

En el corredor diagonal, los conjuntos finitos de prefijos forman un sistema
inverso $X_{m+1}\to X_m$. El espacio de historias infinitas es
$X_\infty=\varprojlim X_m$, mientras sus álgebras de funciones forman el
sistema inductivo opuesto. Esta dualidad de direcciones ---inversa para
historias, directa para observables--- será el soporte exacto de la nueva
bigraduación.

\subsection{Doble proyección y metrología}

HMT no identifica el número real publicado con el estado que lo genera. Una
familia compatible de cilindros produce una lectura arquimediana; otra
proyección organiza la salida excepcional y sus simetrías. La realidad
numérica aparece como sombra de una doble proyección de un estado más rico.
Los teoremas de generación monodrómica de $\pi$, $e$, $\varphi$ y $\alpha$,
los lectores de acción y
gravitación, la ley de masas, los sectores de mezcla y violación CP, y la
genealogía excepcional se aplican en sus dominios declarados. El puente hace
explícitos los objetos operatoriales que transportan sus conclusiones.

La identidad de Planck usa el mismo bloque $B_c$, su
transporte en $SO^+(1,1)$ y la identidad autodual de la recta de masas. El
hecho matemático decisivo para esta traducción es que «observable físico» y
«historia generadora» son tipos diferentes. Una aplicación UCP puede publicar
el primero; la coordenada $m$ y la categoría TPK conservan la segunda.

\begin{sourcebox}
Las premisas del puente son la inclusión unital
$a\mapsto a\otimes I_9$, el bloque $B_c=9P_++5P_-$, la identidad autodual de
la recta de masas y la periodicidad de orden \(108\). Cada una se utiliza con
el dominio y las hipótesis de su teorema; las conclusiones matriciales se
deducen mediante los mapas escritos en este capítulo.
\end{sourcebox}


% El capítulo se distribuye en dos dependencias causales. La unión de los dos
% selectores reproduce todas sus líneas científicas exactamente una vez.
\ifdefined\HMTGMCContinuumMPS
\section{Límites genealógicos, canales de transferencia y memoria de Stinespring}
\label{p12-83-c22-s7:gmc:chap:continuo-mps}

El continuo HMT y el límite continuo de estados de producto matricial (MPS)
responden a preguntas diferentes. El primero prolonga historias finitas y
después las lee; el segundo refina una red tensorial y exige que su canal de
transferencia admita raíces compatibles. La comparación es fértil sólo si se
construye el canal que los conecta.

\subsection{Separación tipada del límite genealógico y el límite MPS}

En el corredor HMT, una historia infinita es un elemento de
$X_\infty=\varprojlim X_m$. Bajo las condiciones de finitud, sobreyectividad,
compactos anidados, contracción y cobertura coherente demostradas en el
anexo del continuo, el cociente por la lectura arquimediana reconstruye el
compacto real declarado. El estado completo y su valor real siguen siendo
objetos distintos.

Para una MPS translacionalmente invariante, un tensor
$A=\{A^i\in M_D\}_{i=1}^d$ determina el canal de transferencia
\[
 E_A=\sum_i A^i\otimes\overline{A^i}.
 \tag{GMC-7.1}
\]
De las Cuevas, Schuch, Pérez-García y Cirac demuestran que una MPS en forma
irreducible tiene límite continuo si y sólo si $E_A$ es infinitamente
divisible; equivalentemente, existen un canal proyector $P$ y un generador
de Lindblad $L$ con
\[
 E_A=Pe^L,
 \qquad P^2=P,
 \qquad PLP=PL.
 \tag{GMC-7.2}
\]
Tiene límite continuo grueso si alguna potencia $E_A^r$ es infinitamente
divisible \citep{delascuevas2018continuum}.

\begin{controlbox}
La existencia del límite inverso HMT no implica \textup{(GMC-7.2)}. Para aplicar el
teorema MPS hay que exhibir un tensor local uniforme, identificar su canal de
transferencia y comprobar divisibilidad infinita. Recíprocamente,
\textup{(GMC-7.2)} no conserva por sí sola hoja, acarreo, ruta o monodromía.
\end{controlbox}
\fi

\ifdefined\HMTGMCOperatorialTower
\subsection{Refinamiento nonádico mediante esperanza condicional}
\label{p12-83-c22-s7:gmc:sec:refinamiento-condicional}

En la torre
$\mathcal A_{m+1}=\mathcal A_m\otimes M_9$, la esperanza condicional
\[
 \mathbb E_m=\id\otimes\tau_9:
 \mathcal A_{m+1}\longrightarrow\mathcal A_m,
 \qquad \tau_9(b)=\frac1{9}\Tr(b),
 \tag{GMC-7.3}
\]
es UCP y satisface
\[
 \mathbb E_m(a\otimes I_9)=a.
 \tag{GMC-7.4}
\]
En la diagonal $D_{m+1}=C(X_m\times\mathbb Z_9)$ se reduce al promedio
uniforme sobre el último dígito. La evaluación de rama
\[
 \kappa_{m,a}(f)(x)=f(x,a)
 \tag{GMC-7.5}
\]
es otra retracción, esta vez multiplicativa en la diagonal. Promediar todas
las ramas y seleccionar una rama son operaciones distintas, paralelas a las
inclusiones $a\otimes I_9$ y $a\otimes p_a$ del cierre operatorial.

La esperanza \textup{(GMC-7.3)} posee una dilatación con un registro explícito. Sean
\[
 W_a\psi=\frac13\,\psi\otimes e_a,
 \qquad a=0,\ldots,8.
 \]
Entonces
\[
 \mathbb E_m(X)=\sum_{a=0}^8W_a^*XW_a.
 \tag{GMC-7.6}
\]
En forma de Stinespring,
\[
 \pi(X)=X\otimes I_9,
 \qquad
 V\psi=\frac13\sum_{a=0}^8(\psi\otimes e_a)\otimes e_a,
 \qquad
 \mathbb E_m(X)=V^*\pi(X)V.
 \tag{GMC-7.7}
\]
El sistema auxiliar (\emph{ancilla}) final etiqueta exactamente el dígito descartado por el promedio.

\begin{bridgebox}
La fórmula \textup{(GMC-7.7)} es una memoria nonádica canónica en el sentido de
Stinespring: lo que se pierde al aplicar $\mathbb E_m$ permanece ortogonalmente
registrado. No se identifica todavía con la memoria TPK. Esa identificación
requiere un entrelazador que transporte también ruta, hoja y dinámica.
\end{bridgebox}

\subsection{Memoria de Stinespring y crecimiento no uniforme}
\label{p12-83-c22-s7:gmc:sec:memoria-no-uniforme}

El \Cref{p12-83-c22-s4:gmc:thm:dilatacion-uhf} da una representación común, en general
infinito-dimensional, para toda familia compatible. No puede exigirse en
general un sistema auxiliar finito fijo que purifique todos los niveles. En efecto, la
familia compatible de trazas normalizadas en $M_{9^m}$ tiene rango $9^m$ en
el nivel $m$, y cualquier purificación del estado correspondiente necesita
una dimensión auxiliar al menos $9^m$. Por tanto,
\[
 \text{compatibilidad UCP en toda profundidad}
 \quad\centernot\Longrightarrow\quad
 \text{memoria auxiliar finita uniforme}.
 \tag{GMC-7.8}
\]
El crecimiento de memoria no es un defecto de la construcción: cuantifica la
cantidad de genealogía que se exige conservar.

La fidelidad dinámica de esta memoria y su comparación con el límite de
estados de producto matricial se desarrollan en
\cref{p12-83-c22-s7:gmc:sec:fidelidad-dinamica}, sin repetir la torre UHF ni la dilatación
de Stinespring aquí demostradas.
\fi

\ifdefined\HMTGMCContinuumMPS
\subsection{Ecuación de fidelidad dinámica}
\label{p12-83-c22-s7:gmc:sec:fidelidad-dinamica}

La esperanza traza-preservante \(\id\otimes\tau_9\), su dilatación de
Stinespring y la imposibilidad general de una memoria auxiliar finita
uniforme han quedado construidas en
\cref{p12-83-c22-s7:gmc:sec:refinamiento-condicional,p12-83-c22-s7:gmc:sec:memoria-no-uniforme}. Sobre
esa premisa operatorial, y sin abreviar sus pruebas, se formula ahora la
condición adicional de fidelidad genealógica.

Sea $T_m^{\rm TPK}$ el transporte entre registros a profundidad $m$, y sea
$J_m$ la codificación de esos registros en el espacio de dilatación. La
memoria de Stinespring es TPK-fiel si existe un transporte
$\widetilde T_m$ en el ambiente tal que
\[
 J_{m+1}T_m^{\rm TPK}=\widetilde T_mJ_m
 \tag{GMC-7.9}
\]
y las compresiones observacionales satisfacen la naturalidad correspondiente.
Si dos historias distintas se identifican bajo $J_m$, la dilatación puede ser
operatorialmente válida pero no genealógicamente fiel.

La ecuación \textup{(GMC-7.9)} es falsable: basta una pareja de rutas que TPK
distinga y el ambiente identifique, o una composición cuya imagen no coincida
con la composición de las imágenes.

\subsection{Criterio de compatibilidad con límites de estados de producto matricial}

Para aplicar \textup{(GMC-7.2)} al continuo HMT deben completarse cinco pasos:

\begin{enumerate}
 \item escoger una celda local uniforme y construir matrices
 $A^i\in M_D$;
 \item demostrar que el bloqueo nonádico de celdas entrelaza la evolución
 con una isometría de refinamiento;
 \item calcular el canal $E_A$ y su forma irreducible;
 \item resolver si $E_A$ o alguna potencia $E_A^r$ es infinitamente
 divisible;
 \item transportar al límite la hoja, el cociclo de acarreo y la
 monodromía, en vez de conservar sólo el estado tensorial.
\end{enumerate}

Si el canal resultara unitario, $E_A=\Ad_U$, tendría raíces unitarias de todo
orden por cálculo funcional. Pero aún habría que demostrar que ese canal es
el canal de transferencia asociado a un tensor HMT y que sus raíces respetan el refinamiento
genealógico. Si apareciera una parte proyectiva $P$, la memoria HMT ofrecería
una interpretación concreta de los sectores que sobreviven al agrupamiento
de escala.

\subsection{Descomposiciones tensoriales con simetría}

El trabajo de De las Cuevas con Netzer y colaboradores sobre
descomposiciones tensoriales aproximadas y sobre complejos simpliciales
permite otra entrada \citep{delascuevas2021approximate,
delascuevas2024simplicial}. APP aporta un complejo finito con acción de
traslaciones; TPK aporta etiquetas y compatibilidades sobre rutas. Una
descomposición tensorial HMT debería declarar:

\begin{itemize}
 \item qué tensor codifica una celda o una transición;
 \item qué grupo preserva las etiquetas y qué parte sólo actúa tras
 reducción;
 \item si la factorización es exacta, aproximada o sólo asintótica;
 \item qué rango se mantiene estable al aumentar $m$.
\end{itemize}

La oportunidad bidireccional es concreta. Su formalismo aporta criterios de
factorización, invariancia y límite; HMT aporta un sistema de refinamiento con
fibras, acarreo y monodromía que puede producir ejemplos donde una
factorización visible existe pero ninguna factorización genealógicamente fiel
conserva la ruta.

\begin{workbox}
Un artículo posterior puede estudiar «divisibilidad infinita con memoria de
raíces»: no sólo si $E=E_k^k$, sino si se puede elegir una torre coherente de
raíces $E_{k\ell}^\ell=E_k$ equipada con los registros TPK. Esta condición es
más fuerte que la divisibilidad del canal desnudo y constituye una extensión
natural del problema MPS desde HMT.
\end{workbox}
\fi


\section{Realización matricial completamente positiva de la fibración nonádica de observables}
\label{p12-83-c22-s8:gmc:chap:simplectica}

La geometría simpléctica no forma parte de los teoremas publicados sobre
cuadrados mágicos de De las Cuevas--Netzer. El vínculo legítimo pasa por un
corredor vecino y estándar: operadores de Weyl, espacio de fases discreto y
bases mutuamente insesgadas \citep{gibbons2004phase,planat2007rings}. En HMT
ese corredor adquiere una elevación nonádica y una memoria de fibra.

\subsection{La forma alternante y la conmutación de Weyl}

Sobre $V_3=\mathbb F_3^2$ definimos
\[
 \langle(q,p),(q',p')\rangle_{\rm sp}=qp'-pq'.
 \tag{GMC-5.1}
\]
Si $X|k\rangle=|k+1\rangle$,
$Z|k\rangle=\omega^k|k\rangle$ y $\omega=e^{2\pi i/3}$, los operadores
$W_{q,p}=X^qZ^p$ satisfacen, con estas convenciones,
\[
 W_uW_v=\omega^{-\langle u,v\rangle_{\rm sp}}W_vW_u.
 \tag{GMC-5.2}
\]
Por tanto, dos operadores conmutan exactamente cuando sus etiquetas son
ortogonales para la forma alternante. En dimensión dos, las subrectas
isotrópicas maximales son los cuatro puntos de
\[
 \mathbb P^1(\mathbb F_3)=\{[1:0],[0:1],[1:1],[1:2]\}.
 \tag{GMC-5.3}
\]
Cada una genera un contexto abeliano maximal de $M_3(\mathbb C)$; sus
proyectores espectrales constituyen una \PVM. Contextos asociados a líneas
distintas producen las cuatro bases MUB de \textup{(GMC-2.13)}.

\begin{proposition}[Descomposición qutrit por direcciones]
\label{p12-83-c22-s8:gmc:prop:descomposicion-masa-qutrit}
Sean $\mathcal D_a$ las cuatro subálgebras diagonales determinadas por las
líneas de \textup{(GMC-5.3)} y $\mathcal D_a^0$ sus partes de traza cero. Entonces
\[
 M_3(\mathbb C)
 =\mathbb CI\oplus
   \bigoplus_{a\in\mathbb P^1(\mathbb F_3)}\mathcal D_a^0
 \tag{GMC-5.4}
\]
como suma ortogonal para el producto de Hilbert--Schmidt. La cuenta de
dimensiones es $9=1+4\cdot2$.
\end{proposition}

\begin{proof}
Cada $\mathcal D_a^0$ tiene dimensión dos. El insesgamiento implica
$\Tr(A^*B)=0$ para $A\in\mathcal D_a^0$ y
$B\in\mathcal D_b^0$, $a\ne b$. Todas son ortogonales a $I$. Las
dimensiones suman nueve, la dimensión de $M_3(\mathbb C)$.
\end{proof}

Esta identidad explica por qué las doce proyecciones no son doce piezas
independientes arbitrarias: son cuatro resoluciones de la unidad cuyas partes
centradas descomponen todo el espacio operatorial qutrit.

\subsection{La elevación \texorpdfstring{$12\to4\times3$}{12 a 4 por 3}}

Sobre el módulo nonádico $V_9=(\mathbb Z/9\mathbb Z)^2$ se conserva la forma
alternante
\[
 \langle(a,b),(c,d)\rangle_9=ad-bc\pmod9.
 \tag{GMC-5.5}
\]
La línea proyectiva $\mathbb P^1(\mathbb Z/9\mathbb Z)$ está formada por
los vectores primitivos, módulo multiplicación por una unidad. Puede
representarse como
\[
 \{[1:t]:t\in\mathbb Z/9\mathbb Z\}
 \ \sqcup\ 
 \{[3u:1]:u\in\mathbb Z/3\mathbb Z\},
 \tag{GMC-5.6}
\]
y tiene $9+3=12$ puntos. La reducción módulo tres da el diagrama
\[
 \mathbb P^1(\mathbb Z/9\mathbb Z)
 \xrightarrow{\ r\ }
 \mathbb P^1(\mathbb F_3),
 \qquad |r^{-1}(a)|=3.
 \tag{GMC-5.7}
\]

La base $a$ de un contexto qutrit y el resultado $k$ de su \PVM también
forman un conjunto $4\times3$. Una \emph{calibración} es una biyección
\[
 \Xi:\mathbb P^1(\mathbb Z/9\mathbb Z)
 \longrightarrow
 \{P_{a,k}:a\in\mathbb P^1(\mathbb F_3),\ k\in\mathbb F_3\}
 \tag{GMC-5.8}
\]
que envía cada fibra de $r$ a los tres resultados del contexto situado sobre
$a$.

\begin{controlbox}
La fibración es canónica; la asignación del orden interno de cada fibra a los
tres resultados $k$ es un calibre. Una biyección no es todavía una
equivarianza. Para promover $\Xi$ a entrelazador debe demostrarse que las
acciones lineales, proyectivas y antiunitarias pertinentes se transportan en
todas las fibras. El calibre preliminar falla una comprobación antiunitaria en una
fibra; ese fallo se conserva como falsador, no se oculta.
\end{controlbox}

\subsection{Profundidad, acarreo y memoria sobre el espacio de fases discreto}

La geometría de Weyl sobre $\mathbb F_3^2$ explica conmutación, contextos y
MUB. HMT añade tres niveles que no se recuperan de la forma simpléctica por
sí sola:

\begin{enumerate}
 \item la elevación de una dirección ternaria a tres puntos de
 $\mathbb P^1(\mathbb Z/9\mathbb Z)$;
 \item la hoja, el acarreo y la orientación que distinguen estados con la
 misma reducción;
 \item la ruta TPK que transporta esas coordenadas y puede regresar a la
 fase sin regresar al estado.
\end{enumerate}

La aportación potencial a la literatura qutrit no es «otra construcción de
cuatro MUB», ya conocida. Es una geometría de \emph{raíces de contexto}: cada
contexto ternario posee tres antecedentes nonádicos, y esos antecedentes
pueden llevar una dinámica y una monodromía. El cuadrado de orden doce del
\Cref{p12-83-c22-s3:gmc:thm:qms-12-3} es la primera sombra normalizada de ese atlas.

\subsection{Separación tipada entre el solapamiento qutrit
\texorpdfstring{$1/3$}{1/3} y el lector HMT
\texorpdfstring{$1/\pi$}{1/pi}}

En una pareja de bases MUB qutrit,
\[
 \Tr(P_{a,k}P_{b,\ell})=\frac13.
 \tag{GMC-5.9}
\]
En el lector HMT aquí definido, la relación entre una hoja real y una
lectura volumétrica puede adoptar el cociente
\[
 \frac{\pi}{\pi^2}=\frac1\pi.
 \tag{GMC-5.10}
\]
Las dos expresiones no son intercambiables: $1/3$ es un solapamiento entre
proyectores en dimensión tres; $1/\pi$ es un peso escalar de otro lector.
La diferencia no impide coordinarlas. Fijada una \PVM
$\{P_0,P_1,P_2\}$, definimos una medición de nueve efectos:
\[
 F_k^{(\pi)}=\frac3\pi P_k\quad(k=0,1,2),
 \qquad
 F_r^{(\pi)}=\frac{1-3/\pi}{6}I_3\quad(r=3,\ldots,8).
 \tag{GMC-5.11}
\]
Es una \POVM porque $3/\pi<1$ y
\[
 \sum_{r=0}^{8}F_r^{(\pi)}
 =\frac3\pi I_3+\left(1-\frac3\pi\right)I_3=I_3.
 \tag{GMC-5.12}
\]
Si $Q$ es un proyector de rango uno de un contexto MUB distinto,
\[
 \Tr(QF_k^{(\pi)})
 =\frac3\pi\Tr(QP_k)
 =\frac3\pi\cdot\frac13
 =\frac1\pi.
 \tag{GMC-5.13}
\]

\begin{bridgebox}
La fórmula \textup{(GMC-5.11)} transforma exactamente la probabilidad MUB $1/3$ en
el lector $1/\pi$ mediante el calibre $3/\pi$, sin identificar ambos números.
Los seis efectos residuales conservan la unidad y completan una \POVM de
nueve resultados, de la que se obtiene otro \QMS cíclico $(9,3)$. Esta
construcción no identifica todavía la procedencia geométrica de la masa
residual; fija el tipo operatorial exacto que deberá tener.
\end{bridgebox}

\subsection{Grupo simpléctico y ecuación de equivarianza}

En dimensión dos sobre $\mathbb F_3$, el grupo simpléctico coincide con
$SL(2,\mathbb F_3)$ y permuta las cuatro líneas. Su acción se levanta,
proyectivamente, mediante el grupo de Clifford sobre las \PVM qutrit. Sea
$G_{\rm HMT}$ un grupo o grupoide de transformaciones del atlas nonádico y
$U_g$ el operador unitario o antiunitario candidato en la carta qutrit. La
condición que convierte una calibración en un puente dinámico es
\[
 \Xi(g\cdot\ell)=U_g\,\Xi(\ell)\,U_g^*,
 \qquad
 g\in G_{\rm HMT},\quad
 \ell\in\mathbb P^1(\mathbb Z/9\mathbb Z).
\tag{GMC-5.14}
\]
Si $g$ es una flecha y no sólo una simetría global, la ecuación debe incluir
origen y destino y se formula como transformación natural. La verificación de
\textup{(GMC-5.14)} es el paso que separa un etiquetado afortunado de una
transducción estructural.

\begin{sourcebox}
Las identidades simplécticas y MUB se apoyan en la literatura de fase finita
\citep{gibbons2004phase,planat2007rings}. No se atribuyen a los artículos de
cuadrados mágicos de De las Cuevas--Netzer. La fibración nonádica, la
calibración y los lectores de ruta constituyen la contribución propia de la
construcción HMT aquí estudiada.
\end{sourcebox}
